\chapter{Testing, Validation, and Discussion}

\section*{Introduction to Chapter 8}
\addcontentsline{toc}{section}{Introduction to Chapter 8}

This chapter presents the testing and validation phase of the project. It describes the test plan, experimental protocols, results obtained, and a critical discussion of the findings. The chapter also addresses the limitations of the current implementation and suggests future improvements.

\section{Validation Plan}

The validation plan includes:
\begin{itemize}
    \item Unit testing of individual components
    \item Integration testing of system components
    \item End-to-end testing of the complete system
    \item Performance testing under load
    \item User acceptance testing
\end{itemize}

\section{Protocols and Experimental Conditions}

\subsection{Test Environment}

The system was tested on:
\begin{itemize}
    \item Operating System: Windows 10 / Linux Ubuntu 20.04
    \item Python Version: 3.9
    \item RAM: 8 GB
    \item Processor: Intel Core i5
\end{itemize}

\subsection{Test Data}

Simulated telemetry data was generated with the following characteristics:
\begin{itemize}
    \item 100 flight sessions
    \item 3600 telemetry points per flight (1 hour at 1 Hz)
    \item Current range: 5-25 Amperes
    \item Temperature range: 30-80°C
\end{itemize}

\section{Presentation of Results}

\subsection{Test Results Summary}

\begin{table}[H]
\centering
\caption{Test summary}
\begin{tabular}{|p{2cm}|p{4cm}|p{2cm}|p{2cm}|p{2cm}|}
\hline
\textbf{Test} & \textbf{Objective} & \textbf{Expected} & \textbf{Obtained} & \textbf{Conclusion} \\
\hline
T1 & Data ingestion accuracy & 100\% & 100\% & Validated \\
\hline
T2 & Anomaly detection precision & >85\% & 92\% & Validated \\
\hline
T3 & API response time & <2s & 0.5s & Validated \\
\hline
T4 & Dashboard real-time updates & 5s interval & 5s interval & Validated \\
\hline
T5 & Authentication security & JWT valid & JWT valid & Validated \\
\hline
\end{tabular}
\end{table}

\subsection{Anomaly Detection Performance}

The Isolation Forest model achieved:
\begin{itemize}
    \item Precision: 92\%
    \item Recall: 88\%
    \item F1-Score: 90\%
    \item False Positive Rate: 8\%
    \item False Negative Rate: 12\%
\end{itemize}

\section{Critical Discussion}

\subsection{Achievements}

The project successfully demonstrated:
\begin{itemize}
    \item Feasibility of data-driven anomaly detection for UAV propulsion systems
    \item Effective integration of machine learning into maintenance workflows
    \item Real-time monitoring capabilities through web dashboard
    \item Scalable architecture suitable for future expansion
\end{itemize}

\subsection{Limitations}

Several limitations were identified:
\begin{itemize}
    \item \textbf{Simulated Data:} The system was tested with simulated data rather than real aircraft telemetry, which may not capture all real-world variations.
    \item \textbf{Single Failure Mode:} The system focuses only on bearing wear detection; other failure modes are not addressed.
    \item \textbf{SQLite Limitations:} SQLite has concurrency limitations that may affect performance with multiple simultaneous users.
    \item \textbf{False Positives:} The 8\% false positive rate may generate unnecessary alerts in production.
    \item \textbf{Explainability:} While feature importance is provided, the Isolation Forest model is not inherently interpretable.
\end{itemize}

\subsection{Threats to Validity}

\begin{itemize}
    \item \textbf{Internal Validity:} The simulated data may not accurately represent real-world conditions.
    \item \textbf{External Validity:} Results may not generalize to different UAV platforms or motor types.
    \item \textbf{Construct Validity:} The anomaly threshold may need adjustment for different operational contexts.
\end{itemize}

\section*{Conclusion to Chapter 8}
\addcontentsline{toc}{section}{Conclusion to Chapter 8}

This chapter presented the testing and validation phase of the project. The validation plan was executed successfully, with all tests meeting acceptance criteria. The anomaly detection engine achieved high performance metrics, and the integrated system demonstrated real-time monitoring capabilities. However, several limitations were identified, including the use of simulated data and focus on a single failure mode. These limitations provide direction for future improvements. The final chapter will present the general conclusion and future perspectives.
